\section{Phase 6: Basic Properties of BLOM}
\label{sec:phase6_basic_properties}

The purpose of this section is to establish the first layer of closed-form properties of BLOM before attempting any convergence theorem.
The guiding principle is that all statements in this section must be either:
\begin{enumerate}
    \item exact structural consequences of the local analytic map constructed in Phase~4, or
    \item exact consequences of the global assembly mechanisms introduced in Phase~5.
\end{enumerate}
No asymptotic statement is used in this section.

Throughout, we adopt the notation of Phases~0--5.
In particular, the canonical setting is
\[
p:[0,\mathcal T]\to\mathbb R,
\qquad
s=4,\qquad
k=2,
\qquad
\theta=(q_1,\dots,q_{M-1},T_1,\dots,T_M)^\top,
\]
with
\[
q_0,\dots,q_M\in\mathbb R,
\qquad
T_i>0,
\qquad
\mathcal T=\sum_{i=1}^{M}T_i.
\]
For every segment \(i\), the raw BLOM-Analytic coefficient map derived in Phase~4 has the form
\begin{equation}
\label{eq:phase6_local_map}
c_i^{\mathrm{BLOM},2}
=
\mathcal M_{2,i}\bigl(q_{i-2},q_{i-1},q_i,q_{i+1},T_{i-1},T_i,T_{i+1}\bigr),
\end{equation}
for all interior segments \(2\le i\le M-1\), with the obvious truncation near the global boundaries.
The associated local segment polynomial is
\begin{equation}
\label{eq:phase6_segment_poly}
p_i^{\mathrm{loc}}(\tau)
=
\bigl(c_i^{\mathrm{BLOM},2}\bigr)^\top \beta_s(\tau),
\qquad
\tau\in[0,T_i].
\end{equation}

In this section, unless otherwise specified, ``BLOM coefficient'' refers to the raw local map \eqref{eq:phase6_local_map}, i.e.\ the coefficient delivered by the local window centered at segment \(i\).
Assembly-dependent properties will be stated separately.

\subsection{Neighborhood notation}
\label{subsec:phase6_neighborhood_notation}

For the canonical three-segment stencil, define the local waypoint and time neighborhoods of segment \(i\) by
\begin{equation}
\label{eq:phase6_q_neighborhood}
\mathcal Q_i
:=
\{i-2,i-1,i,i+1\}\cap\{0,1,\dots,M\},
\end{equation}
and
\begin{equation}
\label{eq:phase6_T_neighborhood}
\mathcal T_i^{\mathrm{idx}}
:=
\{i-1,i,i+1\}\cap\{1,\dots,M\}.
\end{equation}
Thus \(\mathcal M_{2,i}\) depends only on the waypoint values with indices in \(\mathcal Q_i\) and the durations with indices in \(\mathcal T_i^{\mathrm{idx}}\).

For later use, define the derivative row operator
\[
B^{(\ell)}(\tau):=\frac{d^\ell}{d\tau^\ell}\beta_s(\tau)^\top.
\]

\subsection{Proposition 1: local support}
\label{subsec:phase6_prop1}

We first prove the strongest structural property of the raw BLOM coefficient map: exact finite local support.

\begin{proposition}[Exact local support of the raw BLOM-Analytic map]
\label{prop:phase6_local_support}
Fix an interior segment \(i\in\{2,\dots,M-1\}\) in the canonical setting \(s=4,\ k=2\).
Then the raw BLOM-Analytic coefficient map
\[
c_i^{\mathrm{BLOM},2}
=
\mathcal M_{2,i}\bigl(q_{i-2},q_{i-1},q_i,q_{i+1},T_{i-1},T_i,T_{i+1}\bigr)
\]
satisfies
\begin{equation}
\label{eq:phase6_dcdq_zero}
\frac{\partial c_i^{\mathrm{BLOM},2}}{\partial q_j}=0,
\qquad
j\notin \mathcal Q_i,
\end{equation}
and
\begin{equation}
\label{eq:phase6_dcdT_zero}
\frac{\partial c_i^{\mathrm{BLOM},2}}{\partial T_j}=0,
\qquad
j\notin \mathcal T_i^{\mathrm{idx}},
\end{equation}
whenever the derivatives exist.
Equivalently, the Jacobians
\[
\frac{\partial c^{\mathrm{BLOM},2}}{\partial q},
\qquad
\frac{\partial c^{\mathrm{BLOM},2}}{\partial T}
\]
are exactly banded in the segment index.
\end{proposition}

\begin{proof}
By Phase~4, the raw coefficient vector of segment \(i\) is obtained by the composition
\[
c_i^{\mathrm{BLOM},2}
=
\Lambda_8(T_i)\,C_4^{-1}\,y_i,
\]
where
\[
y_i
=
\begin{bmatrix}
q_{i-1}\\ D(T_i)\chi_{i-1}^\star\\ q_i\\ D(T_i)\chi_i^\star
\end{bmatrix},
\qquad
x_i^{\mathrm{loc},\star}
=
\begin{bmatrix}
\chi_{i-1}^\star\\ \chi_i^\star
\end{bmatrix}
=
\bigl(A_2^{(i)}(T)\bigr)^{-1}B_2^{(i)}(q,T).
\]
In the canonical \(k=2\) construction, the local system matrices satisfy
\[
A_2^{(i)}(T)=A_2^{(i)}(T_{i-1},T_i,T_{i+1}),
\]
and
\[
B_2^{(i)}(q,T)=B_2^{(i)}(q_{i-2},q_{i-1},q_i,q_{i+1},T_{i-1},T_i,T_{i+1}).
\]
Therefore the local solution \(x_i^{\mathrm{loc},\star}\), hence \(y_i\), and hence \(c_i^{\mathrm{BLOM},2}\), depend only on the variables
\[
(q_{i-2},q_{i-1},q_i,q_{i+1},T_{i-1},T_i,T_{i+1}).
\]
Consequently, for every \(q_j\) with \(j\notin\mathcal Q_i\), the map \(\mathcal M_{2,i}\) is constant in \(q_j\), so
\[
\frac{\partial c_i^{\mathrm{BLOM},2}}{\partial q_j}=0.
\]
Likewise, for every \(T_j\) with \(j\notin\mathcal T_i^{\mathrm{idx}}\), the map is constant in \(T_j\), so
\[
\frac{\partial c_i^{\mathrm{BLOM},2}}{\partial T_j}=0.
\]
This proves \eqref{eq:phase6_dcdq_zero}--\eqref{eq:phase6_dcdT_zero}.
\end{proof}

\begin{remark}[Scope of Proposition~\ref{prop:phase6_local_support}]
\label{rem:phase6_prop1_scope}
Proposition~\ref{prop:phase6_local_support} is a property of the \emph{raw local coefficient map}, equivalently of Scheme~C before any additional global assembly mechanism is imposed.
It does \emph{not} hold, in general, for the assembled coefficients of Scheme~A or Scheme~B, because both of those schemes introduce global or quasi-global coupling through shared knot states or consensus projection.
Thus exact finite local support is strongest for raw BLOM, and is partially traded away in favor of stronger global continuity in Schemes~A and~B.
\end{remark}

\begin{remark}[General \texorpdfstring{$k$}{k} statement]
\label{rem:phase6_prop1_general_k}
A corresponding exact-support proposition for general \(k\) should read
\[
\frac{\partial c_i^{\mathrm{BLOM},k}}{\partial q_j}=0
\quad\text{and}\quad
\frac{\partial c_i^{\mathrm{BLOM},k}}{\partial T_j}=0
\]
outside the \(k\)-window neighborhood.
However, this general statement requires the existence of a fully explicit analytic map \(\mathcal M_{k,i}\) for arbitrary \(k\), which has not yet been derived in this thesis.
Accordingly, the proof is completed here only for the canonical analytic map \(k=2\).
\end{remark}

\subsection{Proposition 2: interpolation}
\label{subsec:phase6_prop2}

The next property is exact interpolation of the waypoint data.

\begin{proposition}[Exact interpolation]
\label{prop:phase6_interpolation}
Let \(p\) be any BLOM trajectory generated by Scheme~A, Scheme~B, or Scheme~C in the canonical setting.
Then
\begin{equation}
\label{eq:phase6_interpolation}
p(t_i)=q_i,
\qquad i=0,1,\dots,M.
\end{equation}
\end{proposition}

\begin{proof}
We argue segment by segment.

For the raw local map, the coefficient vector of segment \(i\) is reconstructed by the Hermite formula
\[
c_i
=
\mathcal H_i(q_{i-1},q_i,T_i,\eta_{i-1},\eta_i),
\]
where the endpoint values \(q_{i-1}\) and \(q_i\) are explicitly inserted into the reconstruction data.
By the unisolvence of degree-\((2s-1)\) Hermite interpolation, the resulting segment polynomial satisfies
\[
p_i(0)=q_{i-1},
\qquad
p_i(T_i)=q_i.
\]
This holds irrespective of whether the lower-order endpoint jets \(\eta_{i-1},\eta_i\) come from Scheme~A, Scheme~B, or are implicit in Scheme~C.

Therefore, for every interior knot \(t_i\),
\[
p_i(T_i)=q_i
\qquad\text{and}\qquad
p_{i+1}(0)=q_i.
\]
Hence the assembled trajectory takes the exact waypoint value \(q_i\) at knot \(t_i\).
The boundary values \(p(0)=q_0\) and \(p(\mathcal T)=q_M\) follow identically from the first and last segments.
\end{proof}

\begin{corollary}[Universal \(C^0\) property]
\label{cor:phase6_C0}
Every BLOM trajectory generated by any of the three assembly schemes belongs to \(C^0([0,\mathcal T])\).
\end{corollary}

\begin{proof}
Equation \eqref{eq:phase6_interpolation} implies that the left and right traces at every interior knot both equal \(q_i\), so the value jump vanishes.
\end{proof}

\subsection{Proposition 3: post-assembly continuity}
\label{subsec:phase6_prop3}

The continuity statement must be formulated at the level of the \emph{final assembled trajectory}.
It is not enough to say that the local window solution is \(C^{2s-2}\) within its own window, because that property does not automatically survive global extraction and assembly.

We therefore state the result scheme by scheme.

\begin{proposition}[Post-assembly continuity]
\label{prop:phase6_continuity}
Let \(p\) be the final assembled BLOM trajectory.

\begin{enumerate}
    \item \textbf{Scheme A.}
    If \(p=p^A\) is assembled by shared junction states as in Phase~5, then
    \begin{equation}
    \label{eq:phase6_schemeA_cont}
    \llbracket p^{(\ell)}\rrbracket(t_i)=0,
    \qquad
    \ell=0,1,\dots,s-1,
    \qquad
    i=1,\dots,M-1.
    \end{equation}
    Hence
    \[
    p^A\in C^{s-1}([0,\mathcal T]).
    \]

    \item \textbf{Scheme B.}
    If \(p=p^B\) is assembled by overlapping consensus as in Phase~5, then
    \begin{equation}
    \label{eq:phase6_schemeB_cont}
    \llbracket p^{(\ell)}\rrbracket(t_i)=0,
    \qquad
    \ell=0,1,\dots,s-1,
    \qquad
    i=1,\dots,M-1.
    \end{equation}
    Hence
    \[
    p^B\in C^{s-1}([0,\mathcal T]).
    \]

    \item \textbf{Scheme C.}
    If \(p=p^C\) is assembled by raw central-segment extraction, then
    \begin{equation}
    \label{eq:phase6_schemeC_C0}
    \llbracket p\rrbracket(t_i)=0,
    \qquad i=1,\dots,M-1.
    \end{equation}
    Therefore
    \[
    p^C\in C^0([0,\mathcal T]).
    \]
    Moreover, for any fixed \(r\in\{1,\dots,s-1\}\), if the endpoint-jet compatibility condition
    \begin{equation}
    \label{eq:phase6_schemeC_compatibility}
    \Gamma_i^{-}(q,T)=\Gamma_i^{+}(q,T)
    \quad\text{up to order }r
    \end{equation}
    holds for all interior knots \(t_i\), then
    \[
    p^C\in C^r([0,\mathcal T]).
    \]
\end{enumerate}
\end{proposition}

\begin{proof}
Parts (1) and (2) follow from the corresponding assembly theorems of Phase~5.
Indeed, Scheme~A and Scheme~B reconstruct both neighboring segments using the same shared or consensus knot-state \(\eta_i\) at the common knot \(t_i\), and Proposition~\ref{prop:phase5_basic_hermite} implies continuity up to order \(s-1\).

Part (3) follows from Corollary~\ref{cor:phase6_C0}, which gives exact continuity of the value.
The conditional \(C^r\) statement is exactly the sufficient condition proved for Scheme~C in Phase~5: if the lower-order endpoint jets match up to order \(r\), then the corresponding derivative jumps vanish up to order \(r\).
\end{proof}

\begin{remark}[What Proposition~\ref{prop:phase6_continuity} does not claim]
\label{rem:phase6_no_C2sminus2}
Proposition~\ref{prop:phase6_continuity} deliberately does \emph{not} claim global \(C^{2s-2}\) continuity for any scheme.
The stronger within-window continuity from Phase~3 is a \emph{local optimality property}, whereas the present proposition is a \emph{global assembly property}.
Without an additional theorem linking overlapping windows at higher derivative orders, these two notions must be kept distinct.
\end{remark}

\subsection{Proposition 4: weak time sensitivity}
\label{subsec:phase6_prop4}

We now study the weakest time-sensitivity statement that is both useful and rigorously provable at the current stage.
Rather than claiming monotonicity, which is generally too strong, we only prove:

\begin{enumerate}
    \item the local coefficient map is continuously differentiable with respect to the local duration vector;
    \item its Jacobians with respect to local durations are uniformly bounded on compact admissible subsets.
\end{enumerate}

As in Proposition~\ref{prop:phase6_local_support}, this proposition is proved for the canonical explicit analytic map of Phase~4.

\paragraph{Local duration domain.}
For an interior segment \(i\in\{2,\dots,M-1\}\), define the local duration vector
\[
\tau_i:=(T_{i-1},T_i,T_{i+1})^\top\in\mathbb R_{>0}^3.
\]
Let
\begin{equation}
\label{eq:phase6_domain}
\mathcal D_i
:=
\left\{
(q_{i-2},q_{i-1},q_i,q_{i+1},\tau_i)
\in \mathbb R^4\times \mathbb R_{>0}^3
\ \middle|\
A_2^{(i)}(\tau_i)\ \text{is nonsingular}
\right\}.
\end{equation}
By Phase~4, on the canonical admissible BLOM domain the matrix \(A_2^{(i)}(\tau_i)\) is symmetric positive definite, hence \(\mathcal D_i\) contains all physically admissible local data.

\paragraph{Canonical analytic map.}
Recall from Phase~4 that
\begin{equation}
\label{eq:phase6_x_def}
x_i^{\mathrm{loc},\star}
=
\bigl(A_2^{(i)}(\tau_i)\bigr)^{-1} B_2^{(i)}(q,\tau_i),
\end{equation}
and
\begin{equation}
\label{eq:phase6_ci_def}
c_i^{\mathrm{BLOM},2}
=
\Lambda_8(T_i)\,C_4^{-1}\,y_i,
\end{equation}
where
\begin{equation}
\label{eq:phase6_y_def}
y_i
=
\begin{bmatrix}
q_{i-1}\\
D(T_i)\chi_{i-1}^\star\\
q_i\\
D(T_i)\chi_i^\star
\end{bmatrix},
\qquad
x_i^{\mathrm{loc},\star}
=
\begin{bmatrix}
\chi_{i-1}^\star\\
\chi_i^\star
\end{bmatrix}.
\end{equation}

\begin{proposition}[Weak time sensitivity of the canonical BLOM-Analytic map]
\label{prop:phase6_time_sensitivity}
Consider the canonical BLOM-Analytic map \(c_i^{\mathrm{BLOM},2}\) defined by
\eqref{eq:phase6_x_def}--\eqref{eq:phase6_y_def} for an interior segment \(i\).
Then:

\begin{enumerate}
    \item \textbf{Continuous differentiability.}
    The map
    \[
    (q_{i-2},q_{i-1},q_i,q_{i+1},T_{i-1},T_i,T_{i+1})
    \longmapsto
    c_i^{\mathrm{BLOM},2}
    \]
    is \(C^1\) on \(\mathcal D_i\).

    \item \textbf{Derivative formula for the local jet state.}
    For every \(j\in\{i-1,i,i+1\}\),
    \begin{equation}
    \label{eq:phase6_dx_dT}
    \frac{\partial x_i^{\mathrm{loc},\star}}{\partial T_j}
    =
    \bigl(A_2^{(i)}\bigr)^{-1}
    \left(
    \frac{\partial B_2^{(i)}}{\partial T_j}
    -
    \frac{\partial A_2^{(i)}}{\partial T_j}\,x_i^{\mathrm{loc},\star}
    \right).
    \end{equation}

    \item \textbf{Derivative formula for the segment coefficient.}
    For every \(j\in\{i-1,i,i+1\}\),
    \begin{equation}
    \label{eq:phase6_dc_dT}
    \frac{\partial c_i^{\mathrm{BLOM},2}}{\partial T_j}
    =
    \frac{\partial \Lambda_8(T_i)}{\partial T_j} C_4^{-1} y_i
    +
    \Lambda_8(T_i) C_4^{-1} \frac{\partial y_i}{\partial T_j},
    \end{equation}
    where
    \begin{equation}
    \label{eq:phase6_dy_dT}
    \frac{\partial y_i}{\partial T_j}
    =
    \mathbf 1_{\{j=i\}}
    \begin{bmatrix}
    0\\ D'(T_i)\chi_{i-1}^\star\\ 0\\ D'(T_i)\chi_i^\star
    \end{bmatrix}
    +
    \begin{bmatrix}
    0\\ D(T_i)\dfrac{\partial \chi_{i-1}^\star}{\partial T_j}\\ 0\\ D(T_i)\dfrac{\partial \chi_i^\star}{\partial T_j}
    \end{bmatrix}.
    \end{equation}

    \item \textbf{Uniform boundedness on compact admissible sets.}
    Let \(K_i\subset \mathcal D_i\) be compact.
    Then there exists a constant \(C_{K_i}>0\) such that for all local data in \(K_i\),
    \begin{equation}
    \label{eq:phase6_bounded_dcdT}
    \left\|
    \frac{\partial c_i^{\mathrm{BLOM},2}}{\partial T_j}
    \right\|
    \le C_{K_i},
    \qquad
    j=i-1,i,i+1.
    \end{equation}
\end{enumerate}
Moreover, by Proposition~\ref{prop:phase6_local_support},
\[
\frac{\partial c_i^{\mathrm{BLOM},2}}{\partial T_j}=0
\qquad\text{for all } j\notin\{i-1,i,i+1\}.
\]
\end{proposition}

\begin{proof}
We proceed statement by statement.

\paragraph{Step 1: \(C^1\)-regularity of the local jet state.}
The map \(A_2^{(i)}(\tau_i)\) is built from the matrices \(S_-(h)\), \(S_+(h)\), and \(H_{\mathrm{mid}}(h)\), each of which is a rational or polynomial matrix function of the positive durations \(h=T_{i-1},T_i,T_{i+1}\).
Likewise, \(B_2^{(i)}(q,\tau_i)\) is built from polynomial and negative-power expressions of the local durations and affine expressions in the local waypoint vector.
Hence both \(A_2^{(i)}\) and \(B_2^{(i)}\) are \(C^\infty\) on the domain \(T_{i-1},T_i,T_{i+1}>0\).

On \(\mathcal D_i\), the matrix \(A_2^{(i)}\) is nonsingular.
Since matrix inversion is smooth on the general linear group, the map
\[
(q,\tau_i)\mapsto \bigl(A_2^{(i)}(\tau_i)\bigr)^{-1}
\]
is \(C^\infty\) on \(\mathcal D_i\).
Therefore
\[
x_i^{\mathrm{loc},\star}
=
\bigl(A_2^{(i)}\bigr)^{-1}B_2^{(i)}
\]
is \(C^1\) on \(\mathcal D_i\).

\paragraph{Step 2: derivative of the local jet state.}
Differentiate the identity
\[
A_2^{(i)}(\tau_i)\,x_i^{\mathrm{loc},\star}=B_2^{(i)}(q,\tau_i)
\]
with respect to \(T_j\), where \(j\in\{i-1,i,i+1\}\).
Using the product rule,
\[
\frac{\partial A_2^{(i)}}{\partial T_j}x_i^{\mathrm{loc},\star}
+
A_2^{(i)}\frac{\partial x_i^{\mathrm{loc},\star}}{\partial T_j}
=
\frac{\partial B_2^{(i)}}{\partial T_j}.
\]
Rearranging,
\[
A_2^{(i)}\frac{\partial x_i^{\mathrm{loc},\star}}{\partial T_j}
=
\frac{\partial B_2^{(i)}}{\partial T_j}
-
\frac{\partial A_2^{(i)}}{\partial T_j}x_i^{\mathrm{loc},\star}.
\]
Since \(A_2^{(i)}\) is invertible on \(\mathcal D_i\), multiplying by \(\bigl(A_2^{(i)}\bigr)^{-1}\) gives \eqref{eq:phase6_dx_dT}.

\paragraph{Step 3: derivative of the coefficient vector.}
Differentiate
\[
c_i^{\mathrm{BLOM},2}
=
\Lambda_8(T_i)\,C_4^{-1}\,y_i.
\]
Because \(C_4^{-1}\) is constant,
\[
\frac{\partial c_i^{\mathrm{BLOM},2}}{\partial T_j}
=
\frac{\partial \Lambda_8(T_i)}{\partial T_j}C_4^{-1}y_i
+
\Lambda_8(T_i)C_4^{-1}\frac{\partial y_i}{\partial T_j}.
\]
This is \eqref{eq:phase6_dc_dT}.
To compute \(\partial y_i/\partial T_j\), note that
\[
y_i
=
\begin{bmatrix}
q_{i-1}\\
D(T_i)\chi_{i-1}^\star\\
q_i\\
D(T_i)\chi_i^\star
\end{bmatrix}.
\]
Only the scaling matrix \(D(T_i)\) depends explicitly on \(T_i\), while \(\chi_{i-1}^\star\) and \(\chi_i^\star\) depend on all three local durations through \(x_i^{\mathrm{loc},\star}\).
Hence
\[
\frac{\partial}{\partial T_j}\bigl(D(T_i)\chi_{i-1}^\star\bigr)
=
\mathbf 1_{\{j=i\}}D'(T_i)\chi_{i-1}^\star
+
D(T_i)\frac{\partial \chi_{i-1}^\star}{\partial T_j},
\]
and similarly for the right-end jet.
This yields \eqref{eq:phase6_dy_dT}.

\paragraph{Step 4: boundedness on compact admissible sets.}
Let \(K_i\subset\mathcal D_i\) be compact.
Because all functions
\[
A_2^{(i)},\ B_2^{(i)},\ \partial_{T_j}A_2^{(i)},\ \partial_{T_j}B_2^{(i)},\ 
\Lambda_8,\ \partial_{T_i}\Lambda_8,\ D,\ D'
\]
are continuous on \(\mathcal D_i\), and because inversion is continuous on the nonsingular set, all functions
\[
x_i^{\mathrm{loc},\star},\quad
\frac{\partial x_i^{\mathrm{loc},\star}}{\partial T_j},\quad
c_i^{\mathrm{BLOM},2},\quad
\frac{\partial c_i^{\mathrm{BLOM},2}}{\partial T_j}
\]
are continuous on \(K_i\).
Every continuous function on a compact set is bounded.
Hence there exists \(C_{K_i}>0\) such that \eqref{eq:phase6_bounded_dcdT} holds for all data in \(K_i\).

Finally, the zero derivative outside the local support follows from Proposition~\ref{prop:phase6_local_support}.
\end{proof}

\begin{remark}[Marked limitation for Proposition~\ref{prop:phase6_time_sensitivity}]
\label{rem:phase6_prop4_limitation}
Proposition~\ref{prop:phase6_time_sensitivity} is completed only for the canonical explicit BLOM-Analytic map derived in Phase~4, namely \(s=4,\ k=2\).
A general \(C^1\)-regularity theorem for arbitrary \((s,k)\) would require an explicit analytic local system
\[
A_k^{(i)}(T)x_i^{\mathrm{loc}}=B_k^{(i)}(q,T)
\]
for all \(k\), together with a corresponding general Hermite reconstruction map.
That general proof is \textbf{not completed in this chapter}.
\end{remark}

\subsection{Consequences for Jacobian structure}
\label{subsec:phase6_jacobian_structure}

Combining Propositions~\ref{prop:phase6_local_support} and \ref{prop:phase6_time_sensitivity}, we obtain the exact sparsity pattern of the canonical BLOM-Analytic Jacobian.

\begin{corollary}[Exact Jacobian bandedness of raw BLOM]
\label{cor:phase6_jacobian_banded}
For the canonical \(s=4,\ k=2\) raw BLOM-Analytic map,
\[
c_i^{\mathrm{BLOM},2}
=
\mathcal M_{2,i}(q_{i-2},q_{i-1},q_i,q_{i+1},T_{i-1},T_i,T_{i+1}),
\]
one has
\[
\frac{\partial c_i^{\mathrm{BLOM},2}}{\partial q_j}=0
\quad\text{for } j\notin\{i-2,i-1,i,i+1\},
\]
and
\[
\frac{\partial c_i^{\mathrm{BLOM},2}}{\partial T_j}=0
\quad\text{for } j\notin\{i-1,i,i+1\}.
\]
Therefore the raw Jacobians with respect to \(q\) and \(T\) are exactly banded.
\end{corollary}

\begin{proof}
This is merely a restatement of Proposition~\ref{prop:phase6_local_support}.
\end{proof}

\subsection{What has and has not been established in Phase~6}
\label{subsec:phase6_summary}

We summarize the exact status of the four basic propositions.

\begin{enumerate}
    \item \textbf{Proposition 1 (local support): completed.}  
    It is rigorously true for the raw BLOM-Analytic coefficient map, equivalently for Scheme~C before any global coupling mechanism is added.

    \item \textbf{Proposition 2 (interpolation): completed.}  
    It is rigorously true for all assembly schemes.

    \item \textbf{Proposition 3 (post-assembly continuity): completed.}  
    The exact statement depends on the scheme:
    \begin{itemize}
        \item Scheme~A and Scheme~B: \(C^{s-1}\);
        \item Scheme~C: \(C^0\) always, and \(C^r\) only under explicit compatibility conditions.
    \end{itemize}

    \item \textbf{Proposition 4 (weak time sensitivity): completed only in the canonical analytic case \(s=4,\ k=2\).}  
    A full general-\((s,k)\) theorem is not yet available.
\end{enumerate}

Thus Phase~6 establishes a complete first layer of BLOM properties without invoking any asymptotic approximation theorem:
\[
\boxed{
\text{local support}
\;\Longrightarrow\;
\text{interpolation}
\;\Longrightarrow\;
\text{scheme-dependent assembly continuity}
\;\Longrightarrow\;
\text{weak time sensitivity}.
}
\]
These are precisely the ``hard'' properties that should be secured before attempting the convergence theory of later phases.